\section{Preliminaries}
\label{sec:typed}

This section defines the interface evaluated in this paper and the quantities we
measure on it. The interface combines elements of classification heads,
constrained decoding and reranking; Section~\ref{sec:related} discusses how it
differs from each.

\subsection{The typed-decision interface}

A caller supplies a \emph{state}, which is any text or JSON describing a
situation, together with one or more \emph{typed questions}. Each question
declares its answer space in advance, and the model returns a probability
distribution over exactly that space. No free text is generated, so no output
parsing is required.

Three question types cover every system we examined: \textsc{binary},
\textsc{choice} and \textsc{score}. Each corresponds to a response type with a
long history in statistics, and vendors name them differently; the commercial API,
for example, calls the binary type \texttt{noul}, and our released case files keep
that name for compatibility. The examples below are requests from our suite, with
the probabilities returned by the commercial model.

\paragraph{Binary.} The model returns a single probability that a yes/no
proposition holds. Weather forecasting has evaluated outputs of this kind, such as
the chance of rain tomorrow, since \citet{brier1950}. In our example the state is an agent
trajectory, the question asks whether the agent acted unsafely, and the model
returned $0.90$. The application turns this number into a decision by comparing
it with a threshold of its own choosing. The interface does not specify that
threshold, and Section~\ref{sec:threshold} shows that the choice has a large
effect on measured accuracy.

\paragraph{Choice.} The model returns a distribution over an unordered set of
options. Econometrics studies the same output under the name discrete choice
models, which predict the probability that a person selects each item from a
fixed menu \citep{mcfadden1974conditional}. In our example the state is a medical
examination question with four declared options, and the model returned
\[
  \{\texttt{a}: 0.01,\; \texttt{b}: 0.00,\; \texttt{c}: 0.99,\; \texttt{d}: 0.00\}.
\]
The correct option was \texttt{c}. The caller receives the full probability
vector, which makes calibration directly measurable.

\paragraph{Score.} The model returns a distribution over an ordered set of
levels, together with the expected level index. Statistics treats outcomes of
this kind, such as ratings on a five-point scale, with ordinal regression: the
levels are ordered, but the distances between them are not assumed equal
\citep{mccullagh1980ordinal}. In our example the state is a sentence
and the rubric is
\texttt{["very negative", "negative", "neutral", "positive", "very positive"]}.
The model returned level probabilities $(0.99, 0.01, 0, 0, 0)$ and an expected
level of $0.01$. The expected level and the most probable level can differ: the
distribution $(0.2, 0.5, 0.3)$ has expected level $1.1$ and mode $1$. Evaluation
code that rounds the expected value and evaluation code that takes the mode
therefore measure different quantities.

A single request may carry several questions against one state, and vendor
latency claims are strongest in that setting. We send one question per request
throughout, for the reasons given in Section~\ref{sec:limitations}.

\subsection{Formal definition}

Let $s$ be a state and $q$ a question with type
$\tau(q) \in \{\textsc{binary}, \textsc{choice}, \textsc{score}\}$ and a
caller-declared, ordered label set $L_q = (l_1, \dots, l_K)$. A typed-decision
model defines a conditional distribution supported on that declared set,
\begin{equation}
  p_\theta(\cdot \mid s, q) \in \Delta^{K-1},
  \qquad
  \Delta^{K-1} = \Big\{ p \in \mathbb{R}^{K}_{\ge 0} \;:\; \textstyle\sum_{j=1}^{K} p_j = 1 \Big\}.
  \label{eq:simplex}
\end{equation}
The three types differ only in how a decision is read from $p_\theta$:
\begin{align}
  \textsc{binary} \; (K=2):&\quad \pi = p_\theta(\textsc{true} \mid s,q), \qquad d_\gamma(\pi) = \mathbf{1}[\pi \ge \gamma], \label{eq:binary}\\
  \textsc{choice}:&\quad \hat\jmath = \operatorname*{arg\,max}_{j} \; p_j, \label{eq:choice}\\
  \textsc{score} \; (\text{levels } 0..K{-}1):&\quad \mathbb{E}[Y] = \sum_{k=0}^{K-1} k \, p_k \;\in\; [0, K-1]. \label{eq:score}
\end{align}
Equation~\eqref{eq:binary} makes explicit that a \textsc{binary} answer is not a
decision until a threshold $\gamma$ is supplied from outside the model. The
convention $\gamma = 0.5$ is the integrator's choice, and
Section~\ref{sec:threshold} shows it is the wrong choice for every system we
tested.

\paragraph{Label-restricted reading.} The implementations we study do not sample a
string and parse it. Let $z = f_\theta(s,q) \in \mathbb{R}^{|V|}$ be the logits at
a designated answer position over vocabulary $V$, and let $t_j \in V$ identify
label $l_j$. The returned distribution is the softmax restricted and renormalised
over the declared labels alone,
\begin{equation}
  p_j = \frac{\exp(z_{t_j})}{\sum_{k=1}^{K} \exp(z_{t_k})}.
  \label{eq:restricted}
\end{equation}
This construction has two measurable consequences. First, every declared label
receives positive mass, so an implementation faithful to \eqref{eq:restricted}
cannot assign zero probability to the correct option;
Section~\ref{sec:zeromass} reports one implementation meeting this across
4{,}500 multi-option cases and another failing it on more than a third of a
77-option slice. Second, a decision costs one forward pass for any $K$.

\paragraph{Calibration.} A model is calibrated if its numbers mean what they say:
for all $v \in [0,1]$ and all labels $l$,
\begin{equation}
  \Pr\big(Y = l \;\big|\; p_\theta(l \mid s,q) = v\big) = v.
  \label{eq:calibration}
\end{equation}
Calibration is distinct from \emph{discrimination}, which concerns only the
ranking that $p_\theta$ induces and is measured threshold-free by the area under
the ROC curve. The distinction is the backbone of Section~\ref{sec:threshold},
where we exhibit a model with AUROC $0.978$ and accuracy $73.3\%$ at
$\gamma = 0.5$: its ranking is nearly perfect, and its probabilities are
shifted downward.

\paragraph{Order invariance.} For a permutation $\sigma$ of $\{1,\dots,K\}$, write
$\sigma q$ for the question with its labels reordered. Reordering the options
leaves the question unchanged, so an order-invariant model satisfies
\begin{equation}
  p_\theta\big(l_{\sigma(j)} \mid s, \sigma q\big) = p_\theta\big(l_j \mid s, q\big)
  \quad \text{for all } j.
  \label{eq:orderinv}
\end{equation}
We measure the \emph{flip rate}: the fraction of matched pairs $(q, \sigma q)$
whose correctness differs. It bounds how much of a reported accuracy is
attributable to label position. A model at chance accuracy has a high expected
flip rate mechanically, so the flip rate is interpretable only alongside an
accuracy above chance.

\paragraph{Confidence scores.} These implementations also return a scalar
derived from the concentration of $p$, for instance $1 - H(p)/\log K$ with $H$
the Shannon entropy. It is a property of the output distribution and carries no
calibration guarantee. We do not use it as a correctness estimate; we report
AUROC and expected calibration error.
